\documentclass[11pt,letterpaper]{article}
\usepackage[margin=1in]{geometry}
\usepackage[T1]{fontenc}
\usepackage{lmodern,microtype,amsmath,amssymb,amsthm,hyperref}
\hypersetup{colorlinks=true,linkcolor=blue,citecolor=blue,urlcolor=blue,
pdftitle={A telescoping verification of a focal-inversion area invariant},
pdfauthor={Alec Kriebel}}
\setlength{\emergencystretch}{2em}
\DeclareMathOperator{\sn}{sn}
\DeclareMathOperator{\cn}{cn}
\DeclareMathOperator{\dn}{dn}
\newtheorem{theorem}{Theorem}
\newtheorem{lemma}[theorem]{Lemma}
\title{A telescoping verification of a\\focal-inversion area invariant}
\author{Alec Kriebel\thanks{ORCID:
\href{https://orcid.org/0009-0001-9320-500X}{0009-0001-9320-500X}.}}
\date{16 October 2026}
\begin{document}
\maketitle
\begin{abstract}
For a confocal elliptical billiard of least period divisible by four,
the signed area of an orbit times that of the polygon obtained by
unit inversion of its vertices about either original focus is constant.
We verify this theorem, already stated and proved by
Ferudun, through an explicit local Jacobi identity: derivative terms
telescope, and a classical cyclic sum reduces the product to a constant.
The calculation gives an explicit value and covers coprime star
windings. An exact triangular family shows why a repeated list whose
length is divisible by four does not satisfy the same assertion.
This is an attributed verification note; theorem novelty and discovery
priority are not claimed. AI tools were used extensively, and the note
has not undergone conventional human peer review.
\end{abstract}

\section{Statement, source and attribution}
For an ordered closed polygon $W=(W_0,\ldots,W_{N-1})$, use its
algebraic area in traversal order,
\[
 A(W)=\frac12\sum_{j=0}^{N-1}\det(W_j,W_{j+1}),\qquad W_N=W_0.
\]
This convention includes self-intersected polygons and counts
repeated traversals with multiplicity. If $F$ is a focus, let
\[
 I_F(X)=F+\frac{X-F}{|X-F|^2}.
\]
The inverse polygon joins the inverted \emph{vertices} by straight
segments. It is not the curved image of the edges.

The observation being checked is $k_{804,a}$ in Table~9, printed
page~11, of the specific version arXiv:2004.12497v11
\cite{reznik}. Garcia and Reznik later list it as $k_{805,a}$
and prove its simple four-periodic value $4$
\cite[Table~1 and Proposition~4.9]{garcia}.
The original table does not explicitly define the period as a least
period. We state that hypothesis precisely and also treat the
stronger interpretation admitting repeated shorter cycles.

\noindent\begin{minipage}{\linewidth}
\begin{theorem}[Known result; verified below]\label{thm}
Let $E$ and $C$ be strictly nested, nondegenerate confocal ellipses,
with $E$ the billiard boundary and $C$ its fixed caustic. Fix a
billiard family of genuine least collision period $N\equiv0\pmod4$.
Then
\[
                 A(P)\,A(I_F(P))
\]
is a positive constant as the initial phase varies. The constant is
the same for the two original foci. Every admitted coprime winding,
including stars, is covered. Reversal preserves the product; $r$
traversals of an admitted primitive cycle multiply it by $r^2$.
\end{theorem}
\end{minipage}\par\medskip

Ferudun's Theorem~1 \cite{ferudun} has this exact scope, including
either focus, stars and repetitions; his Section~6.1 and accompanying
archive already exclude the unrestricted repeated-list assertion.
His proof cancels all possible poles of the area product on a Jacobi
torus. We give a direct telescoping calculation and a compact
boundary fixture. The local partial-fraction and cyclic-sum
mechanisms are classical: see Khare, Lakshminarayan and Sukhatme
\cite[Eq.~(42)]{local} and \cite[Eq.~(20)]{cyclic}, and the generalized
Landen identity of Khare and Sukhatme \cite[Eq.~(13)]{landen}.
We claim no novelty for these mechanisms or the theorem.
Chavez-Caliz's Theorem~6 \cite{chavez} establishes the related
even-period product of a Poncelet polygon's area and the area of
its outer-tangent polygon. The calculation below supplies the
focal-inverse area reduction explicitly.

\section{Canonical variables and area formulas}
Write the squared-axis difference as $c^2$:
\[
 E:\frac{x^2}{a^2}+\frac{y^2}{b^2}=1,\quad
 C:\frac{x^2}{\alpha^2}+\frac{y^2}{\beta^2}=1,\quad
 a>b>0,\quad 0<\beta<b,\quad
 a^2-\alpha^2=b^2-\beta^2>0.
\]
Thus $c^2=a^2-b^2=\alpha^2-\beta^2$. Put
\[
 k=c/\alpha,\quad k'=\beta/\alpha,\quad K=K(k).
\]
Here $k$ is the Jacobi \emph{modulus}; libraries using the parameter
require $k^2$. In the standard confocal parameter
\cite[Theorem~4.3 and Eq.~(4.9)]{stachel},
\begin{equation}\label{canonical}
 P(w)=(-a\sn w,b\cn w),\quad
 v=\frac{2K\tau}{N},\quad \delta=2v,\quad
 a=\frac{\alpha\dn v}{\cn v},\quad b=\frac{\beta}{\cn v}.
\end{equation}
All Jacobi functions below have modulus $k$.
Choose the orientation with $0<\tau<N/2$ and $\gcd(N,\tau)=1$;
then $0<v<K$. The vertices are $P(w+j\delta)$.
The real parametrization has least period $4K$, so this coprimality
condition expresses least collision period, not just closure after
$N$ listed steps. Set
\[
 m=N/2,\qquad \ell=2K/m,\qquad
 S(w)=\sum_{j=0}^{m-1}\dn(w+j\delta).
\]
Since $N$ is divisible by four, $m$ is even and $\tau$ is odd.
Consequently $P(w+m\delta)=-P(w)$, and the sum defining $S$ is a
permutation of the uniform grid of step $\ell$ modulo $2K$.

For shorter formulas put
\[
 h=\sn v,\quad t=h^2,\quad q=1-t,\quad
 U=1-2k^2t+k^2t^2,\quad V=1-2t+k^2t^2,\quad W=1-k^2t^2,
\]
\[
 C_A=\frac{ab\sn v\cn v}{\dn v},\qquad
 C_0=\frac{b\sn v\dn v\,\cn^4v}{\alpha^3}.
\]
The addition laws \cite[Section~22.8]{dlmf} give
\begin{equation}\label{ordinary}
                         A(P_w)=2C_A S(w).
\end{equation}
For clarity, the half-determinant of endpoints
$P(u-v),P(u+v)$ is
\[
 \frac{ab\sn v\cn v\,\dn u}{1-k^2t\sn^2u}
 =\frac{C_A}{2}\{\dn(u-v)+\dn(u+v)\}.
\]
Summing counts each vertex twice, and the $N$-term vertex sum is
twice $S$.

Take $F_+=(c,0)$ and use relative inverse vertices
$Q(w)=(P(w)-F_+)/|P(w)-F_+|^2$; their area equals that of
$I_{F_+}(P_w)$ because the omitted translation is common.
The Euclidean identity
\[
 |P(w)-F_+|^2=(a+c\sn w)^2\ge(a-c)^2>0
\]
also proves that every real inverse vertex is finite. Direct
addition-formula substitution into
$E(u)=\tfrac12\det(Q(u-v),Q(u+v))$ gives
\begin{equation}\label{edge}
 E(u)=\frac{C_0\dn u(1-k^2t\sn^2u)}
 {(1+k\sn u)(U+kV\sn u)^2}.
\end{equation}
For example, before division the determinant of the two focal-relative endpoint
vectors and the product of the \emph{positive focal distances} are respectively
\[
 \frac{2\alpha b\sn v\dn v\,\dn u(1+k\sn u)}
 {1-k^2t\sn^2u},\qquad
 \frac{\alpha^2(1+k\sn u)(U+kV\sn u)}
 {q(1-k^2t\sn^2u)}.
\]
The latter is squared in the inverse determinant denominator.

Pairing antipodal edges, set $B(u)=E(u)+E(u+2K)$ and
$R(u)=\sum_{j=0}^{m-1}B(u+j\delta)$. Then
\begin{equation}\label{pair}
 \frac{B(u)}{2C_0}=
 \frac{(1-k^2t\sn^2u)\{U^2+k^2V(V+2U)\sn^2u\}}
 {\dn u\,(U^2-k^2V^2\sn^2u)^2},
 \qquad A(I_{F_+}(P_w))=R(w+v).
\end{equation}

\section{The local identity and its cyclic sum}
Let $F(u)=1/\dn u$ and define
\[
 \Lambda=
 \frac{U\{k^4t^3+k^4t^2-5k^2t+k^2+2\}}{2k'^4W^3},
 \qquad
 a_0=\frac{q}{(U-V)^2},\qquad a_1=\frac{\Lambda-a_0}{2},
\]
\[
 a_2=\frac{V^2-U^2t}
 {4\sn v\cn v\dn v\,(U^2-k^2V^2)(U-V)}.
\]
The needed identity is
\begin{equation}\label{localidentity}
 \frac{B(u)}{2C_0}
 =a_0F(u)+a_1\{F(u+\delta)+F(u-\delta)\}
       +a_2\{F'(u+\delta)-F'(u-\delta)\}.
\end{equation}
It is a specialization of the classical local partial-fraction
structure \cite{local}; we verify its coefficients algebraically.

Here is a direct check using only addition and differentiation.
If $x=\sn^2u$ and $H=U^2-k^2V^2x$, then
\[
 F(u+\delta)+F(u-\delta)=\frac{2WU\dn u}{H},
\]
\[
 F(u+\delta)-F(u-\delta)
 =\frac{2k^2\sn\delta\cn\delta\,W^2\sn u\cn u}{H}.
\]
Differentiating the last fraction gives, apart from its constant
factor, numerator
\[
 \dn u\{U^2+(k^2V^2-2U^2)x\}.
\]
Use $\sn\delta=2\sn v\cn v\dn v/W$ and $\cn\delta=V/W$.
After multiplication by $\dn u\,H^2$, equation
\eqref{localidentity} is the quadratic polynomial identity
\begin{align*}
 &(1-k^2tx)\{U^2+k^2V(V+2U)x\}\\
 &\quad=a_0H^2+2WU a_1(1-k^2x)H\\
 &\qquad\quad+
 4k^2\sn v\cn v\dn v\,V a_2
 (1-k^2x)\{U^2+(k^2V^2-2U^2)x\}.
\end{align*}
The displayed coefficients match those of $1,x,x^2$, proving
\eqref{localidentity}.

All denominators used here are nonzero on the strict real domain:
\[
 U>0,\quad W>0,\quad U-V=2tk'^2>0,\quad
 U^2-k^2V^2=k'^2W^2>0,\quad \dn u>0.
\]
In particular $U>|kV|$ and $H>0$. There is no division by $V$.
When $V=0$, one has $\delta=K$ and primitivity forces $N=4$.
The coefficient $a_2$ remains finite; its derivative difference
vanishes since $F'$ has period $2K$. Thus this case is included
without a limiting argument.

Sum \eqref{localidentity} over the finite cycle. Since
$m\delta=2K\tau$, the shifted $F$ sums agree and the derivative
difference telescopes to zero. The quarter-shift formula
$\dn(u+K)=k'/\dn u$ therefore gives
\begin{equation}\label{reduction}
                         R(u)=\frac{2C_0\Lambda}{k'}S(u+K).
\end{equation}
Also $\Lambda>0$. Indeed, with $r=k^2$, the bracket in its numerator
is $p(t)=r^2(t^3+t^2)-5rt+r+2$;
$p'(t)=r\{r(3t^2+2t)-5\}<0$ for $0<t<1$, whereas
$p(1)=2(1-r)^2>0$.

\section{A classical grid identity completes the proof}
We include the Fourier verification of the generalized Landen
identity \cite{landen}, so that its normalization is explicit.
Let $\rho=\exp(-\pi K'/K)$, where $K'=K(k')$.
For real $w$, the classical Fourier series
\cite[Eq.~22.11.3]{dlmf} is absolutely and uniformly convergent:
\[
 \dn(w,k)=\frac{\pi}{2K}+
 \frac{2\pi}{K}\sum_{r\ge1}\frac{\rho^r}{1+\rho^{2r}}
                        \cos(r\pi w/K).
\]
A finite uniform-grid sum retains precisely the indices divisible
by $m$. Let $\kappa\in(0,1)$ have nome $\rho^m$, put
$K_m=K(\kappa)$, $\kappa'=\sqrt{1-\kappa^2}$, and
$\lambda=mK_m/K$. The inverse nome formulas
\cite[Section~22.2]{dlmf} or the monotonic real nome map give
existence of $\kappa$. Comparing Fourier coefficients yields
\[
              S(w)=\lambda\dn(\lambda w,\kappa).
\]
Since $\lambda\ell/2=K_m$, the quarter-shift formula gives
\begin{equation}\label{half}
                         S(w)S(w+\ell/2)=\lambda^2\kappa'.
\end{equation}
Finally
\[
 \frac{K+v}{\ell}=\frac{m+\tau}{2}\in\mathbb Z+\tfrac12.
\]
Equations \eqref{ordinary}, \eqref{reduction} and \eqref{half}
prove Theorem~\ref{thm}, with explicit value
\begin{equation}\label{constant}
       A(P_w)A(I_{F_+}(P_w))
        =\frac{4C_A C_0\Lambda}{k'}\,\lambda^2\kappa'>0.
\end{equation}
Central negation is the ordered vertex shift by $m$ and exchanges
the two foci. It simultaneously negates every relative inverse
vertex, preserving its area; thus both focal inverse areas agree.
Reversal negates both areas. A repeated traversal multiplies
each shoelace sum by its repetition count.

For $N=4$, the axial orbit has areas $2ab$ and $2/(ab)$, recovering
the previously proved value $4$ \cite{garcia}. For a circular boundary
of radius $a$ the focus is the center and inversion has vertex radius
$1/a$. Direct regular-star areas give
$N^2\sin^2(2\pi\tau/N)/4$, separately from the noncircular proof.
Hyperbolic and degenerate caustics are outside the theorem.

\section{An exact repeated-list boundary}
The failure of the unrestricted repeated-list interpretation is
already established in \cite{ferudun}. The following different
fixture also supplies an elementary connected-family check.
Take
\[
 E:\frac{x^2}{8}+\frac{y^2}{5}=1,\qquad
 C:\frac{x^2}{32/9}+\frac{y^2}{5/9}=1,\qquad F=(\sqrt3,0),
\]
and the ordered triangles
\[
 P=\bigl((2\sqrt2,0),(-4\sqrt2/3,5/3),(-4\sqrt2/3,-5/3)\bigr),
 \qquad \widehat P=-P.
\]
These are least-period-three billiards in one continuous forward
family. Here is an explicit justification. Scaling
$(x,y)=(2\sqrt2 X,\sqrt5 Y)$ makes the outer ellipse a unit circle
and the inner ellipse have axes $2/3,1/3$.
A chord joining $z=(X,Y)$ and $z'=(U_1,V_1)$ is tangent exactly when
$3XU_1+6YV_1=-2$.
Put $n=(3X,6Y)$, $D=|n|^2\ge9$, and let $J$ rotate by $\pi/2$.
Its two tangent neighbors are
\[
                z_\pm=\frac{-2n\pm\sqrt{D-4}\,Jn}{D}.
\]
They are distinct circle points and neither equals $z$, since
$3X^2+6Y^2>0$. Direct multiplication gives
$3(z_+)_x(z_-)_x+6(z_+)_y(z_-)_y=-2$.
Thus the three tangent chords close for every $z$.
They remain distinct and noncollinear, and their counterclockwise
orientation at $z=(1,0)$ persists continuously.
As $z$ moves to $(-1,0)$ this joins $P$ to $\widehat P$.

Confocality gives the Euclidean reflection law for this tangent
family. In the unscaled ellipse let $N_E=(x/8,y/5)$ and let $v$
be a unit chord direction and put $n_v=Jv$.
Expansion using $x^2/8+y^2/5=1$ gives
\[
 8(n_v)_x^2+5(n_v)_y^2-(n_v\cdot(x,y))^2
                                      =40(N_E\cdot v)^2.
\]
Tangency to $C$ makes the left side $40/9$.
The two inward tangent directions consequently have the same
normal component $N_E\cdot v=-1/3$ and opposite tangential
components, which is precisely specular reflection.
All three contacts lie inside their segments, and three distinct
collision vertices make the least period exactly three.

Direct vertex inversion and shoelace evaluation give
\[
 A(P)=A(\widehat P)=\frac{50\sqrt2}{9},\quad
 A(I_F(P))=\frac{8\sqrt2-\sqrt3}{80},\quad
 A(I_F(\widehat P))=\frac{8\sqrt2+\sqrt3}{80}.
\]
The two products are $(80-5\sqrt6)/72$ and
$(80+5\sqrt6)/72$. Four traversals create lists of length twelve
and multiply each product by sixteen. The resulting difference is
\[
                              -\frac{20\sqrt6}{9}\ne0.
\]
Their least period remains three. This refutes the unrestricted
listed-length assertion and is consistent with Theorem~\ref{thm}.

\section{Dated records and verification limits}
GitHub records PR207 as created on 1 October 2026
\cite{pr207}; the preserved submitted proof is pinned by its commit
identifier in the accompanying chronology. Zenodo's separate
file metadata dates Ferudun's matching PDF and source archive to
2 October 2026, with matching DOI registration that day
\cite{ferudun}. The explicit telescoping reconstruction and the
fixture used here were recorded during our 16 October audit.
These dates describe the observed records. They do not establish
historical firstness, causal independence, or simultaneous discovery,
and later additions are not backdated to the PR.

The accompanying exact checker certifies finite algebraic identities
and the boundary fixture. The geometric checker uses reflected
Euclidean rays and direct focal inversion; its finite numerical
tests are diagnostics, not a proof or an interval certificate.
Neither the proof nor these computations establishes theorem or
method novelty. Numerical evaluation near $k=1$ or $v=0$ can require
substantially increased precision because of cancellation.

AI tools were used extensively in proof development, algebraic
verification, source checking, adversarial review, programming and
drafting. The author is responsible for the note. It is an
unrefereed preprint and has not undergone conventional human
peer review or proof-assistant formalization.

\begin{thebibliography}{10}
\bibitem{ferudun} A.~Ferudun, \emph{Focal Inversion Area Products in
Elliptic Billiards}, 2 October 2026, unrefereed preprint.
\href{https://doi.org/10.5281/zenodo.23093706}{doi:10.5281/zenodo.23093706}.
\bibitem{reznik} D.~Reznik, R.~Garcia and J.~Koiller,
\emph{Eighty New Invariants of $N$-Periodics in the Elliptic Billiard},
\href{https://arxiv.org/abs/2004.12497v11}{arXiv:2004.12497v11} (2020).
\bibitem{garcia} R.~Garcia and D.~Reznik, Exploring self-intersected
$N$-periodics in the elliptic billiard, \emph{Ann. Math. Inform.}
\textbf{55} (2022), 49--75.
\href{https://doi.org/10.33039/ami.2022.02.001}{doi:10.33039/ami.2022.02.001}.
\bibitem{stachel} H.~Stachel, On the motion of billiards in ellipses,
\emph{Eur. J. Math.} \textbf{8} (2022), 1602--1622.
\href{https://doi.org/10.1007/s40879-021-00524-2}{doi:10.1007/s40879-021-00524-2}.
\bibitem{local} A.~Khare, A.~Lakshminarayan and U.~Sukhatme,
\emph{Local Identities Involving Jacobi Elliptic Functions},
\href{https://arxiv.org/abs/math-ph/0306028}{arXiv:math-ph/0306028} (2003).
\bibitem{cyclic} A.~Khare, A.~Lakshminarayan and U.~Sukhatme,
\emph{Cyclic Identities Involving Jacobi Elliptic Functions. II},
\href{https://arxiv.org/abs/math-ph/0207019v2}{arXiv:math-ph/0207019v2} (2003).
\bibitem{landen} A.~Khare and U.~Sukhatme,
\emph{Generalized Landen Transformation Formulas for Jacobi Elliptic Functions},
\href{https://arxiv.org/abs/math-ph/0312074}{arXiv:math-ph/0312074} (2003).
\bibitem{chavez} A.~C.~Chavez-Caliz, More About Areas and Centers of
Poncelet Polygons, \emph{Arnold Math. J.} \textbf{7} (2021), 91--106.
\href{https://doi.org/10.1007/s40598-020-00154-8}{doi:10.1007/s40598-020-00154-8}.
\bibitem{dlmf} NIST Digital Library of Mathematical Functions,
Chapter~22, Sections~22.2, 22.8 and 22.11,
\url{https://dlmf.nist.gov/22} (accessed 16 October 2026).
\bibitem{pr207} A.~Kriebel, GitHub PR207, \emph{5100044: prove the
arXiv k804,a focal-inversion area-product invariant},
\url{https://github.com/AlecKriebel/Math/pull/207}.
\end{thebibliography}
\end{document}
